\section{Conditions Defining Semisimplicity}

\begin{proposition}\label{prop:semisimple-no-unit}
    Let $F$ be a field, let $A$ be an $F$-algebra (not necessarily with identity), and let $V$ be a left $A$-module. Then the following are equivalent:
    \begin{enumerate}
        \item Every submodule $W \subseteq V$ is a direct summand of $V$.
        \item There exists a family $\{V_i\}_{i\in I}$ of irreducible submodules of $V$ such that
              \[
                  V=\bigoplus_{i\in I} V_i.
              \]
    \end{enumerate}
\end{proposition}

\begin{proof}
    We prove $(1)\Rightarrow(2)$ and $(2)\Rightarrow(1)$.

    \medskip
    \noindent\textbf{(1) $\Rightarrow$ (2).}
    Assume that every submodule of $V$ is a direct summand.

    We first show that every nonzero submodule of $V$ contains an irreducible submodule.
    Let $0\neq W\subseteq V$ be a submodule, and choose $0\neq w\in W$.
    Let $P$ be the submodule of $V$ generated by $w$, that is, the smallest submodule containing $w$.
    Since $V$ is an $F$-vector space and every submodule is in particular an $F$-subspace, we have $Fw\subseteq P$, hence $w\in P$ and therefore $P\neq 0$.

    Consider the set of proper submodules of $P$, partially ordered by inclusion.
    Since $\{0\}$ is a proper submodule of $P$, this set is nonempty.
    Moreover, the union of any chain of proper submodules is again a proper submodule of $P$.
    Hence, by Zorn's lemma, there exists a maximal proper submodule $Q\subsetneq P$.

    By assumption, $Q$ is a direct summand of $P$, so there exists a submodule $Q'\subseteq P$ such that
    \[
        P=Q\oplus Q'.
    \]
    We claim that $Q'$ is irreducible.
    Indeed, let $0\neq X\subseteq Q'$ be a submodule.
    Then
    \[
        Q\oplus X
    \]
    is a submodule of $P$ properly containing $Q$.
    By the maximality of $Q$, we must have $Q\oplus X=P$.
    Since also $P=Q\oplus Q'$, it follows that $X=Q'$.
    Thus $Q'$ has no nonzero proper submodule, so $Q'$ is irreducible.

    Therefore every nonzero submodule of $V$ contains an irreducible submodule.

    Now let
    \[
        E:=\sum \{\, S\subseteq V \mid S \text{ is an irreducible submodule of }V \,\}.
    \]
    Then $E$ is a submodule of $V$.
    We claim that $E=V$.
    Suppose not. Then $E$ is a proper submodule of $V$, and by assumption there exists a submodule $U\subseteq V$ such that
    \[
        V=E\oplus U.
    \]
    Since $E\neq V$, we have $U\neq 0$.
    By the first part, $U$ contains an irreducible submodule $S$.
    But then $S\subseteq E$ by the definition of $E$, and also $S\subseteq U$.
    Hence
    \[
        S\subseteq E\cap U=0,
    \]
    a contradiction.
    Thus $E=V$, so $V$ is the sum of irreducible submodules.

    It remains to show that this sum is in fact direct.
    Write
    \[
        V=\sum_{i\in I} V_i
    \]
    with each $V_i$ irreducible.
    Let $J\subseteq I$ be maximal such that
    \[
        \sum_{j\in J} V_j
    \]
    is a direct sum.
    Then
    \[
        \sum_{j\in J} V_j=\bigoplus_{j\in J} V_j.
    \]
    We claim that
    \[
        V=\bigoplus_{j\in J} V_j.
    \]
    Indeed, let $i\in I$.
    Since $V_i$ is irreducible, the intersection
    \[
        V_i\cap \left(\bigoplus_{j\in J} V_j\right)
    \]
    is either $0$ or $V_i$.
    If it were $0$, then
    \[
        \left(\bigoplus_{j\in J} V_j\right)\oplus V_i
    \]
    would still be a direct sum, contradicting the maximality of $J$ unless $i\in J$ already.
    Hence the intersection must be $V_i$, so
    \[
        V_i\subseteq \bigoplus_{j\in J} V_j.
    \]
    Since this holds for every $i\in I$, we get
    \[
        V=\sum_{i\in I} V_i \subseteq \bigoplus_{j\in J} V_j.
    \]
    The reverse inclusion is obvious, so
    \[
        V=\bigoplus_{j\in J} V_j.
    \]
    This proves~(2).

    \medskip
    \noindent\textbf{(2) $\Rightarrow$ (1).}
    Assume that
    \[
        V=\bigoplus_{i\in I} V_i
    \]
    where each $V_i$ is irreducible.
    Let $W\subseteq V$ be a submodule.
    Consider the collection
    \[
        \mathcal{J}:=\left\{\, J\subseteq I \;\middle|\; W\cap \bigoplus_{j\in J} V_j =0 \right\}.
    \]
    By Zorn's lemma, there exists a maximal element $J\in\mathcal{J}$.

    We claim that
    \[
        V=W+\bigoplus_{j\in J} V_j.
    \]
    Suppose not.
    Then there exists some $i\in I$ such that
    \[
        V_i \nsubseteq W+\bigoplus_{j\in J} V_j.
    \]
    Since $V_i$ is irreducible, the intersection
    \[
        V_i\cap \left(W+\bigoplus_{j\in J} V_j\right)
    \]
    is either $0$ or $V_i$.
    The second possibility would imply
    \[
        V_i\subseteq W+\bigoplus_{j\in J} V_j,
    \]
    contrary to our choice of $i$.
    Hence
    \[
        V_i\cap \left(W+\bigoplus_{j\in J} V_j\right)=0.
    \]
    In particular,
    \[
        W\cap \left(\bigoplus_{j\in J} V_j \oplus V_i\right)=0,
    \]
    which contradicts the maximality of $J$.
    Therefore
    \[
        V=W+\bigoplus_{j\in J} V_j.
    \]

    Finally, this sum is direct, because by the definition of $J$ we have
    \[
        W\cap \bigoplus_{j\in J} V_j=0.
    \]
    Thus
    \[
        V=W\oplus \bigoplus_{j\in J} V_j,
    \]
    so $W$ is a direct summand of $V$.
    This proves~(1).
\end{proof}

